\booksection{ملحق الفصل 1: موضوع نظام تشغيل HarmonyOS}

\chapterintro{HarmonyOS هو نظام تشغيلمطورة ذاتياً من Huawei، ومنذ إصداره في 2019، مرّ بعدة إصدارات رئيسية. إنه ليس مجرد “بديل آخر لـ Android”، بل كائن جديد بفلسفة تصميم فريدة. دعونا نتعرف بعمق على نظام التشغيل الصاعد هذا.}

\booksubsection{أولاً: ما هو HarmonyOS؟ ولماذا هو مختلف؟}

فلسفة HarmonyOS الأساسية هي “ترابط كل الأشياء”. على عكس أنظمة التشغيل التقليدية التي تخدم أساساً جهازاً واحداً، صُمم HarmonyOS منذ البداية مع مراعاة التعاون متعدد الأجهزة.

\textbf{البنية الموزعة}هي روح HarmonyOS. ببساطة، يمكنها تحويل أجهزة متعددة إلى “محطة طرفية فائقة” افتراضية. هاتفك، جهازك اللوحي، حاسوبك، ساعتك، تلفازك، لم تعد أفراداً معزولين، بل “أعضاء” في كلٍّ متكامل، يمكنها استدعاء قدرات بعضها البعض.

مثال: أنت تشاهد فيديو على هاتفك وتريد مشاهدته على شاشة كبيرة، الطريقة التقليدية هي العرض على الشاشة؛ أما على HarmonyOS، يمكنك “نقل” الفيديو مباشرة إلى التلفاز لمتابعة المشاهدة، ويتحول الهاتف إلى جهاز تحكم عن بعد. هذا ليس عرضاً بسيطاً على الشاشة، بل التطبيق نفسه يعمل عبر الأجهزة.

\booksubsection{ثانياً: HarmonyOS 4.3: الإصدار الناضج المبني على AOSP}

HarmonyOS 4.3 هو الإصدار المنتشر حالياً على هواتف Huawei. تجدر الإشارة إلى أن HarmonyOS 4.3 متوافق مع إطار كود مشروع Android مفتوح المصدر (AOSP)، لذا يمكنه تشغيل تطبيقات Android بشكل كامل — معظم التطبيقات التي تنزلها من Huawei AppGallery هي في جوهرها تطبيقات Android، تعمل عبر طبقة التوافق في HarmonyOS. لكن HarmonyOS 4.3 ليس مجرد قشرة فوق Android، بل أضاف على أساس AOSP الكثير من القدرات الموزعةمطورة ذاتياً من Huawei، ومحرك Ark وغيرها من التقنيات.

\subsubsection*{الخصائص الأساسية:}

\begin{itemize}
  \item \textbf{توافق تطبيقات Android}: نظراً للتطوير على أساس AOSP، يمكن لـ HarmonyOS 4.3 تشغيل جميع تطبيقات Android تقريباً، والمستخدم لا يكاد يشعر بالفرق.
  \item \textbf{محرك Ark}: مترجم وبيئة تشغيلمطورة ذاتياً من Huawei، يجعل التطبيقات تعمل بسلاسة أكبر وتوفير طاقة أكبر. مقارنة بالإصدارات المبكرة، تحسنت السلاسة بوضوح.
  \item \textbf{الجهاز الطرفي الفائق}: في مركز التحكم، يمكنك رؤية أجهزة HarmonyOS الأخرى المحيطةبشكل بديهي، السحب للاتصال. الهاتف يستدعي لوحة المفاتيح والماوس من الحاسوب، والجهاز اللوحي يستدعي كاميرا الهاتف، كلها مريحة.
  \item \textbf{بطاقات الخدمة}: مشابهة لعناصر iOS المصغّرة، لكن أكثر مرونة. يمكن تحريك أيقونة التطبيق لأعلى لكشف البطاقة، عرض المعلومات الرئيسية، دون فتح التطبيق.
  \item \textbf{الوضع النقي}: يسمح فقط بتثبيت التطبيقات من متجر التطبيقات الرسمي، يقلل من مخاطر البرمجيات الخبيثة.
  \item \textbf{مركز الخصوصية}: إدارة مركزية لأذونات التطبيقات، عرض سلوكيات التطبيقات.
\end{itemize}

\subsubsection*{خصائص واجهة المستخدم:}

\begin{itemize}
  \item \textbf{قفل الشاشة بالمجلة}: تعرض شاشة القفل صوراً جميلة ومقالات، التمرير يدخل إلى سطح المكتب.
  \item \textbf{البطاقات الشاملة}: أيقونات سطح المكتب يمكن توسيعها إلى بطاقات، تعرض معلومات أكثر.
  \item \textbf{محطة التحويل الفائقة}: منطقة تخزين مؤقتة، يمكن تخزين النصوص والصور والملفات مؤقتاً، ثم سحبها إلى أي تطبيق للاستخدام. المشاركة عبر التطبيقات أصبحت سلسة للغاية.
  \item \textbf{البحث الذكي}: بحث شامل، يمكن البحث في التطبيقات، الملفات، الإعدادات، صفحات الويب.
\end{itemize}

\booksubsection{ثالثاً: HarmonyOS NEXT (HarmonyOS 6): نقطة البداية الحقيقية لنظام تشغيل Huawei المستقل}

HarmonyOS NEXT (HarmonyOS NEXT 5.0) يختلف اختلافاً جوهرياً عن HarmonyOS 4.3 — فهو يعتمد نواة HarmonyOS المصغّرةمطورة ذاتياً من Huawei، ولم يعد يعتمد على نواة Linux وكود AOSP، و\textbf{لم يعد يدعم تطبيقات Android مباشرة}. هذه خطوة رئيسية في انتقال HarmonyOS من “شكل انتقالي متوافق مع Android” إلى “نظام تشغيل مستقل تماماً”.

\subsubsection*{التغييرات الجوهرية:}

\begin{itemize}
  \item \textbf{نواة مستقلة بالكامل}: HarmonyOS NEXT يستخدم نواةمطورة ذاتياً من Huawei، لا تحتوي على أي كود AOSP. النظام أكثر إيجازاً، أكثر أماناً، أكثر كفاءة.
  \item \textbf{لم يعد يدعم تطبيقات Android}: هذا أكبر تغيير. HarmonyOS NEXT يمكنه تشغيل تطبيقات HarmonyOS الأصلية فقط، ولا يمكنه تثبيت وتشغيل ملفات APK. قبل الترقية، يرجى التأكد من أن التطبيقات التي تستخدمها عادةً لديها إصدارات HarmonyOS أصلية.
  \item \textbf{النظام البيئي للتطبيقات الأصلية}: Huawei تبذل جهداً كبيراً لدفع المطورين لترحيل تطبيقاتهم إلى منصة HarmonyOS الأصلية. حالياً التطبيقات الرئيسية مثل WeChat و Alipay و Taobao لديها إصدارات HarmonyOS أصلية، لكن بعض التطبيقات الصغيرة قد لا تكون متكيفة بعد.
  \item \textbf{بنية أمان Star Shield}: نظام أمان جديد تماماً، حماية كاملة السلسلة من العتاد إلى البرمجيات. يدعم تحكماً أكثر دقة في الأذونات، مثل "تحديد الموقع التقريبي" — التطبيق يعرف فقط أنك في مكان ما تقريباً، لا يعرف موقعك الدقيق.
  \item \textbf{دمج عميق لـ AI}: قدرات AI على مستوى النظام، ليس فقط تطبيق معين لديه AI. مثلاً ترجمة AI للنصوص يمكنها ترجمة الفيديو في الوقت الفعلي، و AI لقص الصور يمكنه التعرف الذكي على الموضوع.
  \item \textbf{الانتقال السلس عبر الأجهزة}: قدرات موزعة أكثر تقدماً. حالة التطبيق يمكن أن تنتقل بسلاسة بين الأجهزة، المستند الذي تحرره على هاتفك، عند المشي إلى الحاسوب يفتح تلقائياً على الحاسوب.
  \item \textbf{الخدمات الميتة}: خدمات خفيفة يمكن استخدامها دون تثبيت. مثلاً مسح رمز الاستجابة السريعة لفتح وظيفة معينة مباشرة، دون تنزيل التطبيق أولاً.
\end{itemize}

\begin{tcbwarn}{اقرأ قبل الترقية إلى HarmonyOS NEXT}
بعد الترقية إلى HarmonyOS NEXT، لن يتمكن هاتفك من تشغيل تطبيقات Android. يرجى التأكد قبل الترقية: (1) التطبيقات التي تستخدمها عادةً لديها إصدارات HarmonyOS أصلية؛ (2) لا يوجد في عملك أو حياتك سيناريو يعتمد ضرورة على تطبيقات Android. لا يمكن التراجع عن الترقية إلى HarmonyOS 4.3.

\end{tcbwarn}

\subsubsection*{تغييرات منطق التشغيل:}

\begin{itemize}
  \item \textbf{التنقل بالإيماءات}: التمرير من حافة يسرى أو يمنى للرجوع، التمرير لأعلى من الأسفل للعودة إلى سطح المكتب، التمرير لأعلى والتوقف للدخول إلى المهام المتعددة.
  \item \textbf{النوافذ الذكية المتعددة}: يدعم تقسيم الشاشة، النوافذ العائمة، الرؤية المتوازية وغيرها من أشكال المهام المتعددة.
  \item \textbf{مكالمة Changlian}: خدمة مكالمات فيديو من Huawei، تدعم مشاركة الشاشة، التجميل، تمويه الخلفية.
\end{itemize}

\booksubsection{رابعاً: نظام HarmonyOS البيئي: الأجهزة والتطبيقات}

HarmonyOS ليس مجرد نظام للهواتف، بل نظام بيئي يغطي أجهزة متعددة.

\subsubsection*{أنواع الأجهزة المدعومة:}

\begin{itemize}
  \item \textbf{الهواتف والأجهزة اللوحية}: الأجهزة الأساسية، الوظائف الأكثر اكتمالاً.
  \item \textbf{الساعات الذكية والأساور}: تعمل بشكل مستقل أو بالتعاون مع الهاتف.
  \item \textbf{الشاشات الذكية (التلفاز)}: مركز الترفيه المنزلي.
  \item \textbf{مكبرات الصوت الذكية}: مدخل التفاعل الصوتي.
  \item \textbf{السيارات}: قمرة القيادة HarmonyOS، تجربة قمرة القيادة الذكية.
  \item \textbf{المنزل الذكي}: الثلاجات، المكيفات، الغسالات، إلخ.
\end{itemize}

\subsubsection*{الوضع الحالي للنظام البيئي للتطبيقات:}

متجر تطبيقات HarmonyOS يضم ملايين التطبيقات، وتتوفر معظم التطبيقات الرئيسية. تعمل Huawei بنشاط على تشجيع المطورين على إنشاء تطبيقات أصلية لـ HarmonyOS. حتى منتصف 2026، أطلقت WeChat وAlipay وTaobao وDouyin وMeituan وJD.com وغيرها من التطبيقات الرائدة إصدارات أصلية لـ HarmonyOS، والاستخدام اليومي خالٍ من العوائق تقريباً. بعض التطبيقات المتخصصة والتطبيقات القطاعية قد لا تزال قيد التكييف، لذا يُنصح قبل الترقية بالتأكد من أن تطبيقاتك المستخدمة بكثرة مدعومة.

\booksubsection{خامساً: نصائح عملية لـ HarmonyOS}

\begin{itemize}
  \item \textbf{الجهاز الطرفي الفائق}: اسحب من الزاوية اليمنى العلوية للشاشة لأسفل لفتح مركز التحكم، اضغط على "الجهاز الطرفي الفائق"، يمكنك رؤية الأجهزة الأخرى المحيطة، اسحب للاتصال.
  \item \textbf{محطة التحويل الفائقة}: اضغط مطولاً على نص/صورة/ملف، اختر "نسخ إلى محطة التحويل"، ثم في تطبيقات أخرى اضغط مطولاً للصق، أو اسحب للخارج.
  \item \textbf{بطاقات الخدمة}: اسحب أيقونة التطبيق لأعلى لكشف البطاقة، اضغط مطولاً على البطاقة لضبط الحجم أو إضافتها إلى سطح المكتب.
  \item \textbf{النوافذ الذكية المتعددة}: اسحب من حافة الشاشة اليمنى أو اليسرى للداخل وتوقف، افتح الشريط الجانبي، اضغط على التطبيق لفتحه في نافذة عائمة.
  \item \textbf{حماية الخصوصية}: الإعدادات > الخصوصية، يمكنك عرض سجل استخدام أذونات التطبيقات، ضبط تحديد الموقع التقريبي، إلخ.
\end{itemize}

\begin{tcbwarn}{تحذيرات}
نظام HarmonyOS في تحديث وتكرار مستمر، الإصدارات المختلفة قد يكون فيها بعض الاختلافات. هذا الوصف يعتمد على الخصائص الرئيسية لـ HarmonyOS 4.3 و 6، الوظائف المحددة يرجى الرجوع إلى الإصدار الفعلي على هاتفك. يُنصح بالحفاظ على تحديث النظام للحصول على أحدث الوظائف وترقيعات الأمان.

\end{tcbwarn}

\begin{tcbpractice}{جرب الآن: اختبر وظائف HarmonyOS المميزة}
\begin{enumerate}
  \item افتح مركز التحكم، ابحث عن "الجهاز الطرفي الفائق"، انظر إن كانت هناك أجهزة HarmonyOS أخرى حولك.
  \item جرب "محطة التحويل الفائقة": انسخ نصاً إلى محطة التحويل، ثم اسحبه في الملاحظات.
  \item اسحب لأعلى أيقونات بعض التطبيقات الشائعة، انظر ما المحتوى الذي تعرضه بطاقات الخدمة الخاصة بها.
  \item اذهب إلى الإعدادات > الخصوصية، اعرض أي التطبيقات استخدمت أذونات حساسة مؤخراً.
\end{enumerate}

\end{tcbpractice}
